\section*{Chapter \ref{ch:g8:winds-global-circulation} --- Winds and the Global Circulation}

\begin{solution}{exo:g8:winds-global-circulation:1}
(a) The north-east trade wind, from the north-east. (b) The \omterm{def:g8:winds-global-circulation:global-winds}{westerlies},
from the north-west. (c) The \omterm{def:g8:winds-global-circulation:global-winds}{polar easterlies}, from the east (from the
north-east). (d) The south-east trade wind, from the south-east.
\end{solution}

\begin{solution}{exo:g8:winds-global-circulation:2}
(a) Warm air rises and cold air sinks. (b) Low. (c) From high pressure
towards low pressure. (d) To the right.
\end{solution}

\begin{solution}{exo:g8:winds-global-circulation:3}
(a) False: a \omterm{def:g8:winds-global-circulation:breezes}{sea breeze} blows from the \omterm{def:g2:water-on-earth:ocean}{sea} to the land (by day).
(b) True. (c) False: they blow from the east (from the north-east or the
south-east). (d) True.
\end{solution}

\begin{solution}{exo:g8:winds-global-circulation:4}
The drawing shows a loop: along the ground, an arrow from the \omterm{def:g2:water-on-earth:ocean}{sea} to the
land; over the land, an arrow going up; at height, an arrow from the land
back over the \omterm{def:g2:water-on-earth:ocean}{sea}; over the \omterm{def:g2:water-on-earth:ocean}{sea}, an arrow going down. H is marked on the
ground under the sinking air (over the \omterm{def:g2:water-on-earth:ocean}{sea}), L under the rising air (over
the land).
\end{solution}

\begin{solution}{exo:g8:winds-global-circulation:5}
(a) 1.41 times as much ground. (b) 50~\%. (c) 26~\%. (d) Near the
\omterm{def:g3:continents-and-oceans:poles}{poles}, sunlight arrives at a slant and is spread over a large area, so
each square metre receives little energy.
\end{solution}

\begin{solution}{exo:g8:winds-global-circulation:6}
(a) \omterm{def:g8:winds-global-circulation:global-winds}{Polar easterlies} (\ang{70} to \ang{60} N), \omterm{def:g8:winds-global-circulation:global-winds}{westerlies} (\ang{60} to
\ang{30} N), north-east \omterm{def:g8:winds-global-circulation:global-winds}{trade winds} (\ang{30} N to the \omterm{def:g3:continents-and-oceans:equator}{equator}),
south-east \omterm{def:g8:winds-global-circulation:global-winds}{trade winds} (\omterm{def:g3:continents-and-oceans:equator}{equator} to \ang{30} S), \omterm{def:g8:winds-global-circulation:global-winds}{westerlies} (\ang{30} to
\ang{60} S), \omterm{def:g8:winds-global-circulation:global-winds}{polar easterlies} (\ang{60} to \ang{70} S). (b) Near the
\omterm{def:g3:continents-and-oceans:equator}{equator}, in the doldrums. (c) Near \ang{30} N and \ang{30} S, under the
sinking air.
\end{solution}

\begin{solution}{exo:g8:winds-global-circulation:7}
(a) The land ($31 > 22$): the wind blows from the \omterm{def:g2:water-on-earth:ocean}{sea} to the land, a \omterm{def:g8:winds-global-circulation:breezes}{sea
breeze}. (b) The \omterm{def:g2:water-on-earth:ocean}{sea} ($21 > 17$): the wind blows from the land to the
\omterm{def:g2:water-on-earth:ocean}{sea}, a \omterm{def:g8:winds-global-circulation:breezes}{land breeze}. (c) The land: $31 - 17 = 14$ degrees; the \omterm{def:g2:water-on-earth:ocean}{sea}: $22 -
21 = 1$ degree. (d) The Sun's rays heat only a thin layer of ground,
which warms and cools quickly; the \omterm{def:g2:water-on-earth:ocean}{sea} spreads its heat through a thick
layer of water, so its temperature hardly changes between day and night.
\end{solution}

\begin{solution}{exo:g8:winds-global-circulation:8}
(a) From the north. (b) To the right. (c) From the north-east (it is the
north-east trade wind). (d) Without rotation it would blow from the
south; the rotation turns it to the left; it blows from the south-east.
\end{solution}

\begin{solution}{exo:g8:winds-global-circulation:9}
Near the \omterm{def:g3:continents-and-oceans:equator}{equator}, the ground and the \omterm{def:g2:water-on-earth:ocean}{sea} are strongly \emph{heated} by
the Sun; the warm, damp air \emph{rises}; as it rises it \emph{cools},
and its water vapour turns into \emph{\omterm{def:g4:clouds-rain-snow:cloud}{clouds}}, which give showers and
thunderstorms. Near \ang{25} N, the air that rose at the \omterm{def:g3:continents-and-oceans:equator}{equator}
\emph{sinks}; under it the pressure at the ground is \emph{high}; the
sinking air warms, its \omterm{def:g4:clouds-rain-snow:cloud}{clouds} dry up and the region stays \emph{dry}: a
\omterm{def:g5:weather-and-climate:desert}{desert}.
\end{solution}

\begin{solution}{exo:g8:winds-global-circulation:10}
(a) The \omterm{def:g8:winds-global-circulation:global-winds}{westerlies}. No: they blow from the west, against a ship sailing
west. (b) To reach the belt of the \omterm{def:g8:winds-global-circulation:global-winds}{trade winds}, south of about \ang{30}
N. (c) The north-east trade wind, between the \omterm{def:g3:continents-and-oceans:equator}{equator} and about \ang{30}
N. (d) The \omterm{def:g8:winds-global-circulation:global-winds}{westerlies}, north of about \ang{30} N (near \ang{40} N, the
\omterm{def:g6:latitude-longitude:latitude}{latitude} of the home port).
\end{solution}

\begin{solution}{exo:g8:winds-global-circulation:11}
(a) $1 \div 0.5 = 2$: the table gives 2. (b) $1 \div 0.866 \approx
1.1547$, that is \textbf{1.15}. (c) At \ang{30}: $1000 \div 1.15 \approx
\qty{870}{W}$ on each square metre (exactly $1000 \times 0.866 =
\qty{866}{W}$). At \ang{60}: $1000 \div 2 = \qty{500}{W}$.
\end{solution}

\begin{solution}{exo:g8:winds-global-circulation:12}
(a) Three loops: air rises at the \omterm{def:g3:continents-and-oceans:equator}{equator} and sinks near \ang{30};
rises near \ang{60} and sinks near \ang{30}; rises near \ang{60} and sinks
at the \omterm{def:g3:continents-and-oceans:poles}{pole}. (b) L at the \omterm{def:g3:continents-and-oceans:equator}{equator} and near \ang{60}; H near \ang{30} and
at the \omterm{def:g3:continents-and-oceans:poles}{pole}. (c) From \ang{30} to the \omterm{def:g3:continents-and-oceans:equator}{equator}, the surface air goes south
and is turned to the right: north-east trade wind. From \ang{30} to
\ang{60}, it goes north and is turned to the right: it comes from the
south-west. From the \omterm{def:g3:continents-and-oceans:poles}{pole} to \ang{60}, it goes south, turned to the
right: polar easterly. (d) In the middle cell the surface air flows
towards the \omterm{def:g3:continents-and-oceans:poles}{pole}, that is north; turned to the right, it heads towards
the north-east: these are the \omterm{def:g8:winds-global-circulation:global-winds}{westerlies}.
\end{solution}

\begin{solution}{pb:g8:winds-global-circulation:1}
\textbf{1.} $59.6 - 23.5 = 36.1$ degrees of \omterm{def:g6:latitude-longitude:longitude}{longitude}.

\textbf{2.} $36.1 \times 108 \approx 3899$: about \qty{3900}{km}.

\textbf{3.} Yes: $3900 - 3896 = 4$~km, a tiny difference.

\textbf{4.} Yes: both lie near \ang{14} N, between the \omterm{def:g3:continents-and-oceans:equator}{equator} and about
\ang{30} N, the belt of the north-east trade wind.

\textbf{5.} From the north-east (the arrows point towards the
south-west).

\textbf{6.} Forwards, and a little sideways: the wind comes from behind
the ship, slightly from its right, and pushes it towards the west.

\textbf{7.} It is part of the \omterm{def:g8:winds-global-circulation:circulation}{global circulation}: air keeps sinking near
\ang{30} N and rising near the \omterm{def:g3:continents-and-oceans:equator}{equator}, so the surface air keeps flowing
towards the \omterm{def:g3:continents-and-oceans:equator}{equator}, and the Earth's rotation always turns it to the
right.

\textbf{8.} The trade wind would blow against the ship. The captain
should sail north, beyond about \ang{30} N, to find the \omterm{def:g8:winds-global-circulation:global-winds}{westerlies}.

\textbf{9.} $7 \times 24 = \qty{168}{km}$ a day.

\textbf{10.} $3 \times 168 = \qty{504}{km}$.

\textbf{11.} $3896 \div 168 \approx 23.2$ days.

\textbf{12.} About \textbf{23 days}: less than a month.
\end{solution}
